\documentclass[11pt]{article}

% Generic paper layout: 1in margins are a conventional baseline for drafts,
% class papers, and many ad hoc submissions, while still using the page better
% than LaTeX's default article text block.
\usepackage[
    paper=letterpaper,
    margin=1in,
    includeheadfoot,
    headheight=14pt,
    headsep=0.22in,
    footskip=0.35in
]{geometry}

% Typography.
\usepackage[T1]{fontenc}
\usepackage[utf8]{inputenc}
\usepackage{lmodern}
\usepackage{microtype}
\usepackage{textcomp}
\usepackage{setspace}
\setstretch{1.08}
\raggedbottom

% Core math support.
\usepackage{amsmath, amssymb, amsthm, mathtools}
\usepackage{mathrsfs}
\usepackage{dsfont}
\numberwithin{equation}{section}
\allowdisplaybreaks

% Tables, figures, and layout helpers.
\usepackage{graphicx}
\usepackage{booktabs}
\usepackage{threeparttable}
\usepackage{array}
\usepackage{tabularx}
\usepackage{subcaption}
\usepackage{caption}
\usepackage{siunitx}
\usepackage{xcolor}
\usepackage{enumitem}
\usepackage[most]{tcolorbox}
\usepackage{listings}
\graphicspath{{images/}{figures/}}

\tcbset{
  lecrcodeboxstyle/.style={
    colback=gray!8,
    colframe=black!65,
    boxrule=0.45pt,
    arc=2.4mm,
    left=6pt,
    right=6pt,
    top=6pt,
    bottom=6pt
  }
}

\lstdefinelanguage{Rnote}{
  language=R,
  morekeywords={
    mutate,summarise,summarize,filter,group_by,ungroup,across,
    case_when,if_else,coalesce,replace_na,left_join,inner_join,
    full_join,semi_join,anti_join,pivot_longer,pivot_wider,
    ggplot,aes,labs,facet_wrap
  }
}
\lstdefinestyle{notescode}{
  language=Rnote,
  basicstyle=\ttfamily\small,
  keywordstyle=\color{blue!75!black},
  commentstyle=\itshape\color{teal!55!black},
  stringstyle=\color{red!70!black},
  showstringspaces=false,
  columns=fullflexible,
  keepspaces=true,
  breaklines=true,
  breakatwhitespace=false,
  aboveskip=0pt,
  belowskip=0pt,
  frame=none
}
\newtcblisting{rcodebox}{
  enhanced,
  breakable,
  listing only,
  lecrcodeboxstyle,
  listing options={style=notescode}
}

% Citations and links.
\usepackage[authoryear, round]{natbib}
\usepackage{hyperref}
\hypersetup{
    colorlinks=true,
    linkcolor=black,
    citecolor=blue!70!black,
    urlcolor=blue!70!black
}
\usepackage[nameinlink, noabbrev]{cleveref}

% Title and section styling.
\usepackage{titling}
\usepackage{titlesec}
\setlength{\droptitle}{-0.6in}
\pretitle{\begin{center}\LARGE\bfseries}
\posttitle{\par\end{center}\vspace{0.4em}}
\preauthor{\begin{center}\normalsize}
\postauthor{\par\end{center}\vspace{0.15em}}
\predate{\begin{center}\small}
\postdate{\par\end{center}\vspace{1.0em}}

\titleformat{\section}
  {\normalfont\large\bfseries}
  {\thesection}
  {0.75em}
  {}
\titleformat{\subsection}
  {\normalfont\normalsize\bfseries}
  {\thesubsection}
  {0.7em}
  {}
\titleformat{\subsubsection}
  {\normalfont\normalsize\itshape}
  {\thesubsubsection}
  {0.65em}
  {}
\titlespacing*{\section}{0pt}{1.4\baselineskip}{0.6\baselineskip}
\titlespacing*{\subsection}{0pt}{1.0\baselineskip}{0.45\baselineskip}
\titlespacing*{\subsubsection}{0pt}{0.8\baselineskip}{0.35\baselineskip}

\captionsetup{
    font=small,
    labelfont=bf,
    textfont=normalfont,
    justification=centering,
    singlelinecheck=false,
    skip=6pt
}
\setlist{
    itemsep=0.25em,
    topsep=0.45em
}

% Line-break penalties: prevent mid-expression line breaks.
\relpenalty=10000
\binoppenalty=10000

% =====================================================================
% MATH SHORTHAND MACROS
% =====================================================================
\newcommand*\N{\mathbb{N}}
\newcommand*\R{\mathbb{R}}
\newcommand*\Q{\mathbb{Q}}
\newcommand*\Z{\mathbb{Z}}
\newcommand*\E{\mathbb{E}}
\let\oldP\P
\renewcommand*\P{\mathbb{P}}
\let\O\varnothing
\let\oldemptyset\emptyset
\let\emptyset\varnothing
\newcommand{\ind}{\perp\!\!\!\perp}

\DeclareMathOperator*{\argmax}{arg\,max}
\DeclareMathOperator*{\argmin}{arg\,min}
\DeclareMathOperator*{\plim}{plim}
\DeclareMathOperator{\Var}{Var}
\DeclareMathOperator{\Cov}{Cov}
\DeclareMathOperator{\Corr}{Corr}
\DeclareMathOperator{\Bias}{Bias}
\DeclareMathOperator{\MSE}{MSE}
\DeclareMathOperator{\Span}{span}
\DeclareMathOperator{\rank}{rank}
\DeclareMathOperator{\Tr}{Tr}
\DeclareMathOperator{\diag}{diag}
\DeclareMathOperator{\supp}{supp}

\DeclarePairedDelimiter{\ceil}{\lceil}{\rceil}
\DeclarePairedDelimiter{\floor}{\lfloor}{\rfloor}
\DeclarePairedDelimiter{\abs}{\lvert}{\rvert}
\DeclarePairedDelimiter{\norm}{\lVert}{\rVert}

\newcommand{\Prob}{\mathbb{P}}
\newcommand{\Expect}{\mathbb{E}}
\newcommand{\wh}{\widehat}
\newcommand{\wt}{\widetilde}
\newcommand{\hf}{{1/2}}
\newcommand{\lnorm}[2]{\left\lVert #1 \right\rVert_{#2}}
\newcommand{\Norm}[1]{\left\lVert #1 \right\rVert}
\newcommand{\Fnorm}[1]{\lnorm{#1}{\mathrm{F}}}
\newcommand{\fnorm}[1]{\left\lVert #1 \right\rVert_{\mathrm{F}}}
\newcommand{\opnorm}[1]{\left\lVert #1 \right\rVert_{\mathrm{op}}}
\newcommand{\nnorm}[1]{\left\lVert #1 \right\rVert_{*}}
\newcommand{\iprod}[2]{\left\langle #1, #2 \right\rangle}
\newcommand{\pth}[1]{\left( #1 \right)}
\newcommand{\qth}[1]{\left[ #1 \right]}
\newcommand{\sth}[1]{\left\{ #1 \right\}}
\newcommand{\indc}[1]{\mathds{1}_{\left\{ #1 \right\}}}
\newcommand{\calL}{\mathcal{L}}

% =====================================================================
% THEOREM ENVIRONMENTS
% =====================================================================
\theoremstyle{plain}
\newtheorem{theorem}{Theorem}[section]
\newtheorem{lemma}[theorem]{Lemma}
\newtheorem{proposition}[theorem]{Proposition}
\newtheorem{corollary}[theorem]{Corollary}
\newtheorem{assumption}[theorem]{Assumption}
\newtheorem{claim}[theorem]{Claim}

\theoremstyle{definition}
\newtheorem{definition}[theorem]{Definition}
\newtheorem{example}[theorem]{Example}
\newtheorem{notation}[theorem]{Notation}

\theoremstyle{remark}
\newtheorem{remark}[theorem]{Remark}

\crefname{section}{section}{sections}
\Crefname{section}{Section}{Sections}
\crefname{subsection}{subsection}{subsections}
\Crefname{subsection}{Subsection}{Subsections}
\crefname{equation}{equation}{equations}
\Crefname{equation}{Equation}{Equations}
\crefname{theorem}{theorem}{theorems}
\Crefname{theorem}{Theorem}{Theorems}
\crefname{lemma}{lemma}{lemmas}
\Crefname{lemma}{Lemma}{Lemmas}
\crefname{proposition}{proposition}{propositions}
\Crefname{proposition}{Proposition}{Propositions}
\crefname{corollary}{corollary}{corollaries}
\Crefname{corollary}{Corollary}{Corollaries}
\crefname{assumption}{assumption}{assumptions}
\Crefname{assumption}{Assumption}{Assumptions}
\crefname{claim}{claim}{claims}
\Crefname{claim}{Claim}{Claims}
\crefname{definition}{definition}{definitions}
\Crefname{definition}{Definition}{Definitions}
\crefname{example}{example}{examples}
\Crefname{example}{Example}{Examples}
\crefname{remark}{remark}{remarks}
\Crefname{remark}{Remark}{Remarks}

% Notes under tables/figures.
\newcommand{\fignote}[1]{\par\vspace{0.25em}{\footnotesize\textit{Notes:} #1\par}}
\newcommand{\tabnote}[1]{\par\vspace{0.25em}{\footnotesize\textit{Notes:} #1\par}}
